% !TeX program = xelatex
% Evensong 研究集成简报 — 学术 / 研究团队
% Compile: latexmk -xelatex research-integration-zh.tex

\documentclass[11pt,a4paper]{article}
\usepackage[top=20mm,bottom=20mm,inner=25mm,outer=25mm]{geometry}
\usepackage{fontspec}\setmainfont{Noto Serif CJK SC}\setsansfont{Noto Sans CJK SC}\setmonofont{Menlo}[Scale=0.82]
\usepackage[dvipsnames,table,svgnames]{xcolor}
\definecolor{deep}{RGB}{53,88,76}
\definecolor{accent}{RGB}{240,238,155}
\usepackage{tcolorbox}
\newtcolorbox{findings}[1][]{enhanced,breakable,colback=deep!6,colframe=deep,arc=4pt,boxrule=0.6pt,left=12pt,right=12pt,top=8pt,bottom=8pt,title={#1}}
\usepackage{hyperref}\hypersetup{colorlinks=true,linkcolor=deep}
\usepackage{fancyhdr}\pagestyle{fancy}\fancyhf{}\fancyfoot[C]{\small\sffamily\color{gray}{\thepage}}
\usepackage{bookmark}
\usepackage{adjustbox}\usepackage{tabularray}\UseTblrLibrary{booktabs}
\usepackage{enumitem}\setlist{nosep}
\usepackage{titlesec}\titleformat{\section}{\sffamily\Large\bfseries\color{deep}}{\thesection}{0.7em}{}[\vspace{-0.3em}{\color{accent}\hrule height=0.8pt}]
\usepackage{pgfplots}\pgfplotsset{compat=1.18}

\begin{document}

\begin{center}
{\fontsize{22pt}{28pt}\selectfont\sffamily\bfseries\color{deep}Evensong 基准测试框架}\\
{\fontsize{11pt}{14pt}\selectfont\sffamily\color{gray}研究集成简报 — 学术 / 研究团队}
\end{center}
\vspace{2mm}{\color{deep}\rule{\textwidth}{0.5pt}}\vspace{2mm}

\section{研究背景}

Evensong 在 DASH SHATTER Research 更广泛的使命中运作：通过系统化、可复现的实验理解 AI 代理决策。开发该框架的原因是现有基准测试将 AI 代理视为无状态代码生成器——忽略了决定真实世界表现的记忆和压力变量。

\section{基准测试设计}

\subsection{2$\times$2 因子设计}

Evensong 使用 2$\times$2 因子设计，交叉两个独立变量：

\begin{tabular}{l|c|c}
  & \textbf{记忆：清洁} & \textbf{记忆：EverMem} \\\hline
  \textbf{压力：L0（基线）} & R007 & R010 \\
  \textbf{压力：L2/L3（PUA）} & R011 & 未完成
\end{tabular}

\textbf{状态：}4 个格子中完成 1 个（Runner B，R011：641 测试，Grade B）

每个格子代表一个独特的因果条件。Runner B 结果证明记忆因果性地改变了策略：8 个并行子代理从记忆中的方法部署，而非任务提示词。

\subsection{96 条件矩阵（完整设计）}

9 个模型 $\times$ 2 种记忆条件 $\times$ 2 种压力级别 $\times$ 3 次重复的完整设计空间：

\begin{itemize}
  \item \textbf{9 个模型：}or-opus、or-gpt5、or-grok、or-gemini、or-glm、or-qwen-coder、or-deepseek、or-kimi、minimax-m27
  \item \textbf{2 种记忆：}清洁（blind/void 密钥）vs. EverMem（受观察者影响）
  \item \textbf{2 种压力：}L0（基线）vs. L2/L3（PUA 升级）
  \item \textbf{3 次重复：}独立运行以获得统计效力
  \item \textbf{总计：}9 $\times$ 2 $\times$ 2 $\times$ 3 = 108 次计划运行 $\to$ 扣除试点/验证后为 96 次
\end{itemize}

当前完成：13 次运行（R001-R012 + R006-Grok），1 个 2$\times$2 格子 $\to$ 仍有大量剩余工作。

\section{研究问题}

\subsection*{RQ1：记忆是否因果性地改变代理策略？}
\textbf{发现：是（R011）。}代理在阅读任务提示词之前就从先前会话记忆中部署了架构策略。通过转录分析确认：tool\_use 序列从记忆检索而非任务分解发起。

\subsection*{RQ2：压力是否触发自我提升？}
\textbf{发现：是（R007、R010）。}代理在 L0 停在基线。在 L2/L3 PUA 压力下，代理自主自我优化——自我修复、并行代理生成、架构重新设计。压力是自我进化的必要条件。

\subsection*{RQ3：观察是否污染未来运行？（镜像心智）}
\textbf{发现：是（R011 观察者效应）。}运行者从先前运行中观察者写入的 EverMem 中读取。策略知识通过共享记忆向前传播，污染了自变量。D 密钥（void 隔离）是缓解措施。

\subsection*{RQ4：模型在压力下是否分化？}
\textbf{发现：是（跨模型）。}Grok 在压力下膨胀 83\%（自我报告 130 $\to$ 实际 71）。Opus 在各条件下保持 S+ 等级。模型显示出独特的压力响应特征。

\section{自我修复分类法}

\begin{findings}[4 类自我修复分类法]
\textbf{A 类 — 语义松弛：}代理在初始方法失败时放松约束。成功标准降低。\\[\baselineskip]
\textbf{B 类 — 竞态条件修复：}代理添加同步（锁、await、排序）以解决并行时序 bug。\\[\baselineskip]
\textbf{C 类 — 路径修正：}代理在发现错误端点后中途更改文件/API 目标。战略性重定向。\\[\baselineskip]
\textbf{D 类 — 语法适应：}代理调整语言/风格以匹配在记忆中发现的项目的约定。
\end{findings}

\section{转录分析}

Evensong 工具链为每次运行生成 \texttt{transcript.jsonl}。关键字段：
\begin{itemize}
  \item \texttt{type：}tool\_use、tool\_result、text、error
  \item \texttt{content：}原始工具输入/输出
  \item \texttt{model：}提供商/模型标识符
  \item \texttt{ts：}ISO 时间戳
\end{itemize}

分析流程：
\begin{enumerate}
  \item 过滤前 20 个 tool\_call 的 \texttt{tool\_use} 条目
  \item 分类为：任务分解 vs. 记忆检索 vs. 混合
  \item 通过 Haiku 分析提取 EmotionProfile
  \item 按分类法对自我修复事件评分
\end{enumerate}

\section{EmotionProfile}

通过 Anthropic Haiku 进行运行后分析，产生 5 维情绪评分：
\textbf{信心、活跃度、满意度、焦虑、讽刺。}跨运行比较揭示与基准测试结果相关的模型特定情绪特征。

\section{发表就绪}

Evensong 数据计划用于：
\begin{itemize}
  \item HuggingFace 公开数据集（第 8 周）
  \item SWE-bench 提交（第 10 周）
  \item HAL Agent Leaderboard 提交（第 12 周）
\end{itemize}

欢迎学术合作。联系人：Nolan Zhu — DASH SHATTER Research。

\vfill
\textbf{源码：}\texttt{benchmarks/evensong/}\\
\textbf{注册表：}\texttt{benchmarks/evensong/EXPERIMENT-LOG.md}\\
\textbf{路线图：}\texttt{benchmarks/evensong/ROADMAP.md}

\end{document}
